% Section 5, Evaluation (about 3,000 words; docs/paper-mock.md): Tables 1-3 on EMBER at pilot scale. Sources:
% docs/direction-A-A1-design.md (runs 99-100, 105, 108), docs/direction-A-threat-model.md S6 (run 79),
% docs/item-layer-design.md (runs 88-107), docs/direction-A-expert-set.md (run 106), docs/status.md (runs 110, 111).
% A breadth instance, if added, adds rows to Tables 1 and 2 (and possibly a detector table); the domain is
% interchangeable.

We evaluate the three claims of the framework on the EMBER instance: that the side effects it computes are exact
(\Cref{sec:eval:side-effects}), that the directions it derives are more correct than a careful hand analysis and more
complete than constraints learned from attacks (\Cref{sec:eval:directions}), and that detectors built on them resist
every edit in the threat model, including an adaptive attacker who uses the same machinery
(\Cref{sec:eval:detectors}).

\subsection{Setup}
\label{sec:eval:setup}

All numbers are at pilot scale: one seed, a 10\% subsample of each EMBER~2024 split for training and testing, and a
pilot set of 100 files, half malware, for the edits and attacks. Detectors are gradient-boosted
trees~\cite{ke2017lightgbm}, with monotone constraints where a row calls for them. Each detector's threshold is set to
a 1\% false-positive rate on the benign test files, and an attack counts as an evasion when it drives a detected
malware sample below that threshold. Every attacked file is written to disk and extracted again, so no reported score
rests on our own predictions. Unless stated otherwise, the threat model is the full catalogue of
\Cref{tab:instances:catalogue}, with a byte budget of 5\% of the file or 512~KB, whichever is larger.

The attacks come from two sources. The \emph{toolkit} attacks are the eight secml-malware
attacks~\cite{demetrio2021adversarial} that Gimenez et al.\ evaluate against: padding, GAMMA section
injection~\cite{demetrio2021functionality}, GAMMA import injection, two content shifts, DOS-header extension, header
fields, and a format exploit, each with a population of 10 for 10 iterations. We split them into four \emph{seen}
attacks, from which the baseline below is built, and four \emph{held-out} ones. The \emph{adaptive} attack is ours. It
searches the catalogue greedily for the edit that lowers a given detector's score most, scoring each candidate on the
feature vector Theorem~\ref{thm:fw:side-effects} predicts instead of extracting it. It appends bytes and adds
sections, and it aims section names at the hash
buckets the detector relies on (\Cref{sec:fw:directions}). It is a white-box attacker with the same budget of 100
scored candidates as the toolkit attacks, and it writes each kept step to disk with our own byte-level edits.

\subsection{Side Effects Are Exact}
\label{sec:eval:side-effects}

Theorem~\ref{thm:fw:side-effects} says that, under the requirements, the features of an edited file follow from the item
table of the original and the edit. We test it the way Pierazzi et al.\ obtain side effects: apply the edit, extract
the features again, and compare all 2,480 modeled features with the prediction. \Cref{tab:eval:side-effects} reports
the result per operation. A prediction is exact when every feature matches to the extractor's float precision; an edit
is flagged when the engine declines to predict because the edit writes a locating position.

\begin{table}[t]
  \caption{Predicted side effects against re-extraction, on all 2,480 modeled features. Flagged: the engine declined
  to predict (the edit writes a locating position).}
  \label{tab:eval:side-effects}
  \small
  \begin{tabular}{@{}lrrrr@{}}
    \toprule
    Operation & Edits & Exact & Flagged & Wrong \\
    \midrule
    append bytes & 100 & 89 & 11 & 0 \\
    overwrite the DOS header & 100 & 100 & 0 & 0 \\
    random in-place writes & 2,000 & 1,791 & 209 & 0 \\
    add a section & 100 & \synthetic{100} & \synthetic{0} & \synthetic{0} \\
    up to three sections, then append & 1,105 & 1,105 & 0 & 0 \\
    \bottomrule
  \end{tabular}
\end{table}
% SYNTHETIC: the add-section row (exact prediction of the declared delta against re-extraction of our own byte-level
% edit, 100 pilot files). Measured so far: equal to the hand predictor on all 100 files (run 92), and directions
% correct on every edit (run 107). Replace with a C1-exactness run over every declared operation.

No prediction is wrong. The 11 appends the engine flags are exactly the files whose section table points past their
truncated end, where the appended bytes become a section's raw data; the engine finds them by an interval query of the
written positions against the recorded locating positions, with no rule written for this case. The flagged in-place
writes are, likewise, the writes that land on a header field or table entry the parser reads. Combined edits, up to
three added sections followed by an append, are predicted exactly by carrying the item table forward one edit at a
time.

Exactness matters most where hand analysis is weakest: the side effects that are easy to miss. Appending to a file
lowers every section's size ratio, since EMBER divides each section's size by the file's size; it can raise the
largest section entropy, since the overlay counts as a section there; and when the appended bytes follow printable
text, it extends the last string rather than adding one. Each of these follows from the item table without a rule of
its own. The same record supports the other requirement checks, which we report in the appendix: no single-byte change
among 58,666 alters the parse without touching a recorded locating position, and the readouts reproduce the extractor
bit for bit on 2,000 files.

\subsection{Directions: Correct and Complete}
\label{sec:eval:directions}

Before the derivation existed, we wrote the direction table by hand: for each of EMBER's feature groups and each of
the twelve operations, we read the extractor's code and the operation's declared effects and recorded which way the
group's features could move. This is the realistic alternative to the derivation, done by people who knew the
extractor well and had its source open. We then compared it with the table Theorem~\ref{thm:fw:directions} derives.

\begin{table*}[t]
  \caption{Every disagreement between the hand-written direction table and the derived one. In each case the hand rule
  is wrong, and a real file and edit, extracted again, moves the feature in the direction the hand rule excluded. No
  disagreement goes the other way.}
  \label{tab:eval:directions}
  \small
  \begin{tabular}{@{}p{0.17\textwidth}p{0.17\textwidth}lp{0.45\textwidth}@{}}
    \toprule
    Feature & Operation & Hand $\to$ derived & Witness \\
    \midrule
    overlay size & append & up $\to$ both & the appended bytes complete a section that pointed past the truncated end, and
      the overlay shrinks (4 pilot files; e.g.\ 31,482 to 4,096 bytes) \\
    entry section's name & add a section & fixed $\to$ both & on files whose entry point lies in no section, the added
      executable section becomes the entry section (8 pilot files) \\
    Rich header & insert a gap before the PE header & fixed $\to$ both & the parser searches a range that contains the
      gap; content written there creates a Rich header (all 39 pilot files without one) \\
    Rich header & overwrite header fields & fixed $\to$ both & three header fields lie in the same search range (8 of
      the 39) \\
    number of certificates & append & down $\to$ both & completing a signature table past the end adds a certificate (2
      pilot files) \\
    has relocations, exports & append & fixed $\to$ both & completing a directory past the end (relocations: 151 of
      2,000 files; exports: 37) \\
    section entropy & overwrite header fields & fixed $\to$ both & a section whose raw data starts at offset 0 is the
      header (4 of 2,000 files) \\
    \bottomrule
  \end{tabular}
\end{table*}

On every other feature the two tables agree per threat model, and on one feature, the size of the signature table,
the derived direction is tighter than the hand rule. \Cref{tab:eval:directions} lists the disagreements, and each
is the hand rule's error. None of them is exotic. They come from three facts about the parser that are visible only in
its code: it ignores structures that lie past the end of the file, so appending can bring them into existence; it
searches a fixed range for the Rich header, wherever the PE header happens to start; and it picks the entry section by
a rule with a fallback. The derivation needs none of these facts stated: they are locating positions in the item
table, and the interval query does the rest. The first row is the one that would have mattered most, since the overlay
size was constrained as up-only in our first monotone detectors; the adaptive search could not have found the hole,
because its own hand-written predictor shared the same blind spot.

The derived table is sound in the other sense as well. We applied every operation of the catalogue, with adversarial
parameters (repeated, colliding, and case-variant names, empty imports, printable content next to existing strings,
a forged Rich header in the inserted gap), to the 2,000-file set, written by our own byte-level code and extracted
again. Across 73,063 edits, no feature moved in a direction the derived table excludes.

The other realistic alternative is to learn the constraints from attacks. We let each feature move as far as the eight
toolkit attacks moved it, and asked an exact solver for tree ensembles~\cite{kantchelian2016evasion} whether any point
in that region evades each detector. The learned region certifies none of the 47 files the unconstrained detector
catches and none of the 22 the baseline below catches: ranges learned per feature ignore that one edit moves many
features together, and the region they span contains evasions no edit can produce. Learned from the attacks a
defense was built on, the region certifies that defense by construction, including on files our adaptive attack
evades. Computed directions avoid both failures, because they come from what each operation can do rather than from
what a sample of attacks happened to do.

\subsection{Detectors No In-Model Edit Can Evade}
\label{sec:eval:detectors}

Corollary~\ref{cor:fw:robustness} turns the derived table into detectors: drop the features marked both, constrain the rest to
be nondecreasing in the up features and nonincreasing in the down ones, and the score can never fall under any
sequence of edits in the threat model. We train such a detector on the features as they are and on the recorded sums
(\Cref{sec:instances:sums}), and compare it with the alternatives a practitioner would use.

\begin{table*}[t]
  \caption{Detectors under the full threat model. Evasion is over the detected pilot malware. Seen and held-out are
  the toolkit attacks; adaptive is our search. The monotone rows send files outside the population to the
  unconstrained detector.}
  \label{tab:eval:detectors}
  \small
  \begin{tabular}{@{}lrrrrrr@{}}
    \toprule
    Detector & AUC & TPR at 1\% FPR & Detected & Seen & Held-out & Adaptive \\
    \midrule
    unconstrained & 0.996 & 93.5\% & 47/50 & 17\% & 53\% & 8\% \\
    adversarial training on the toolkit & \synthetic{0.995} & \synthetic{92.0\%} & \synthetic{46/50} & \synthetic{2\%} & \synthetic{30\%} & \synthetic{20\%} \\
    feature selection from attacks (PV), 4 attacks & 0.954 & 56.2\% & 22/50 & 0\% & 18\% & 73\% \\
    feature selection from attacks (PV), 8 attacks & 0.944 & 53.8\% & 22/50 & 0\% & --- & \synthetic{77\%} \\
    expert set (40 features), monotone & 0.844 & 29.1\% & 8/50 & 0\% & 0\% & --- \\
    monotone on derived directions, features & 0.992 & 87.2\% & 40/50 & 0\% & 3\%$^\dagger$ & 0\% \\
    monotone on derived directions, sums & 0.994 & 89.6\% & 42/50 & 0\% & 2\%$^\dagger$ & 0\% \\
    \bottomrule
  \end{tabular}
  \\[2pt]
  {\footnotesize $^\dagger$ One header-field evasion, on a file outside the population, scored by the unconstrained
  detector. On files in the population: 0/28 and 0/30.}
\end{table*}
% SYNTHETIC: the adversarial-training row (toolkit adversarial examples added to training; run to be scripted), and
% PV with 8 attacks under the adaptive search after the search fix (run 108 re-ran PV with 4 attacks only; before the
% fix it read 22/22). The expert set has no adaptive cell: the search predicts with the instrumented extractor only.

\Cref{tab:eval:detectors} gives the results. Both monotone rows are evaded by no seen or held-out toolkit attack on
the files the guarantee covers, and by nothing the adaptive search finds: on every file and every candidate edit, the
largest score drop it can produce is exactly 0.0. This is Corollary~\ref{cor:fw:robustness} measured: no edit in the catalogue
lowers a feature the detector reads in the direction it trusts, so the search has nothing to work with. The single
held-out evasion is a header-field attack on a file whose layout is not plain, which the detector routes to the
unconstrained model, where the guarantee does not apply; 38 of the 50 pilot malware, and 1,674 of the 2,000 files,
lie inside the population. Soundness costs a few points of detection: 89.6\% at 1\% false positives on the sums,
against 93.5\% unconstrained. Reading the sums instead of the features recovers about two of those points, because
the re-encodings keep features one-sided that their normalized readouts made two-sided.

The guarantee is per threat model, and the table shows what that means. Under the additions threat model, the monotone
detectors keep more features and detect more (92.2\% on the sums), and the header-field attacks, which lie outside
that threat model, evade them on a third of the files. A derived table is only as wide as the catalogue it was derived
from, and a narrower catalogue buys accuracy at the price of the edits it leaves out.

The closest baseline is Gimenez et al.'s feature selection from perturbation vectors (PV): run attacks on the training
malware, keep the features no attack lowered, and train a monotone model on them. It resists the attacks it was built
from, by construction, and costs over 35 points of detection. It also has a hole that the derivation computes before
any attack runs: 44.4\% of the PV model's gain lies in features the derived table marks as lowerable, most of it in
section-name hash buckets that no toolkit attack happened to lower. The adaptive search fills the hole. It picks
section names whose hash lands in buckets the model trusts, with the sign that lowers them, and evades PV on 16 of the
22 files it detects, using \synthetic{a handful of} extractions per file. Building PV from all eight attacks does not
close the hole: \synthetic{17 of 22}. Of the six files the search does not evade, five have no room for another
section header, and the sixth ends just above the threshold. Adversarial training on the toolkit's outputs keeps the unconstrained detector's accuracy
but, like PV, is robust only to what it was shown: \synthetic{the adaptive search evades it on a fifth of the files}.

The expert set is Gimenez et al.'s 40 hand-picked features, each argued to grow under additions. With the extractor's
two quirks patched it holds under additions, and no toolkit attack evades it, but 40 features keep little signal: it
detects 8 of the 50 pilot malware, at 29.1\% detection. Under the full threat model, 12 of its 40 features can be
lowered by renaming a section, overwriting slack, or rewriting the DOS stub, edits the expert analysis, made for
additions only, does not consider. A hand-picked set inherits the threat model its authors had in mind; a derived one
is recomputed for each catalogue.

Deriving the table costs one pass of the instrumented extractor over the corpus and a sign computation per feature
and operation, \synthetic{about 20 minutes on 2,000 files}. PV needs every attack run against every training malware,
\synthetic{about 40 CPU-hours} at our pilot's budget, and must be repeated for every new attack.
% SYNTHETIC: the cost figures (bookkeeping from the logs of runs 97/107 and 78; not yet collected).
